\chapter{Polynomials and the Derivatives of \texorpdfstring{$\delta$}{delta}}
\label{ch:polynomial}

\begin{quote}
\itshape
If the constant has a spectrum, so does $t$, and so does every polynomial. The answers involve
$\delta'$, $\delta''$, and so on. Before those symbols can mean anything, differentiation has to be
defined for objects that are not functions---and the only definition available is the one that
integration by parts suggests.
\end{quote}

\section{The obstacle, computed}

Take $f(t)=t$. Truncating symmetrically and integrating by parts with $u=t$,
$\dd v=e^{-i\omega t}\dd t$,
\begin{equation}
\int_{-T}^{T}t\,e^{-i\omega t}\dd t
=\left[\frac{t\,e^{-i\omega t}}{-i\omega}\right]_{-T}^{T}
+\frac{1}{i\omega}\int_{-T}^{T}e^{-i\omega t}\dd t
=\frac{2T\cos(\omega T)}{i\omega}+\frac{2\sin(\omega T)}{\omega^{2}} .
\label{eq:t-truncated}
\end{equation}
The first term oscillates with an amplitude that grows linearly in $T$. The failure is therefore
strictly worse than for the constant, where the amplitude was bounded. For $t^{n}$ one gets an
amplitude of order $T^{n}$.

So no amount of truncation will help. The way forward is not analytic but algebraic: obtain the
transform of $t$ from the transform of $1$ by a rule, and postpone the justification of the rule to
Chapter~\ref{ch:tespered}.

\section{The distributional derivative}

\begin{definition}[distributional derivative]\label{def:distderiv}
Let $T$ be a continuous linear functional on $C_{c}^{\infty}(\R)$ (a distribution; see
Chapter~\ref{ch:dist}). Its derivative $T'$ is defined by
\begin{equation}
\ip{T'}{\varphi}:=-\ip{T}{\varphi'},
\qquad
\ip{T^{(n)}}{\varphi}:=(-1)^{n}\ip{T}{\varphi^{(n)}} .
\label{eq:distderiv}
\end{equation}
\end{definition}

The definition is forced by the following observation.

\begin{proposition}[compatibility]\label{prop:compat}
Let $f\in C^{1}(\R)$ with $f'\in\Lonloc$, and let $T_{f}$ denote the distribution
$\varphi\mapsto\int f\varphi$. Then $T_{f}'=T_{f'}$.
\end{proposition}

\begin{proof}
For $\varphi\in C_{c}^{\infty}$, integration by parts on the support of $\varphi$ (which is compact,
so the boundary terms vanish) gives
\begin{equation}
\int_{\R}f'(t)\varphi(t)\dd t=-\int_{\R}f(t)\varphi'(t)\dd t
=-\ip{T_{f}}{\varphi'}=\ip{T_{f}'}{\varphi}. \qedhere
\end{equation}
\end{proof}

Thus \eqref{eq:distderiv} extends the classical derivative wherever the latter exists, and it makes
sense even where no pointwise derivative exists. Two further properties are immediate and will be
used repeatedly: the map $T\mapsto T'$ is linear, and it is continuous with respect to distributional
convergence (Proposition~\ref{prop:derivcont}).

\begin{example}[the derivatives of $\delta$]\label{ex:deltan}
Applying Definition~\ref{def:distderiv} to $\ip{\delta}{\varphi}=\varphi(0)$,
\begin{equation}
\ip{\delta'}{\varphi}=-\ip{\delta}{\varphi'}=-\varphi'(0),
\qquad
\ip{\delta''}{\varphi}=(-1)^{2}\ip{\delta}{\varphi''}=\varphi''(0),
\end{equation}
and in general
\begin{equation}
\ip{\delta^{(n)}}{\varphi}=(-1)^{n}\varphi^{(n)}(0).
\label{eq:deltan}
\end{equation}
So $\delta$ reads off the value at the origin, $\delta'$ reads off (minus) the slope there,
$\delta''$ reads off the curvature. They are different questions about the same point, not
successively sharper peaks.
\end{example}

\begin{remark}[the sign is not decoration]
The factor $(-1)^{n}$ is the boundary term of the integration by parts. Dropping it produces
$\F[t]=-2\pi i\delta'$ instead of $2\pi i\delta'$ and propagates through every formula in this
chapter. It is safer to remember where it comes from than to remember whether it is there.
\end{remark}

\section{Multiplication by \texorpdfstring{$t$}{t} and differentiation in \texorpdfstring{$\omega$}{omega}}

\begin{lemma}[differentiation under the integral]\label{lem:diffint}
Let $f\in\Sw(\R)$. Then $\fh$ is differentiable and
\begin{equation}
\frac{\dd}{\dd\omega}\fh(\omega)=\int_{\R}(-it)f(t)e^{-i\omega t}\dd t .
\label{eq:diffint}
\end{equation}
\end{lemma}

\begin{proof}
Fix $\omega$ and let $h\neq0$. Then
\begin{equation}
\frac{\fh(\omega+h)-\fh(\omega)}{h}
=\int_{\R}f(t)\,\frac{e^{-i(\omega+h)t}-e^{-i\omega t}}{h}\dd t .
\end{equation}
The integrand converges pointwise to $-it\,f(t)e^{-i\omega t}$. To apply dominated convergence,
note that by the mean value theorem the difference quotient is bounded in modulus by
$|t|\,|f(t)|$, which is integrable because $f\in\Sw$. Hence the limit may be taken inside the
integral.
\end{proof}

\begin{theorem}[multiplication by $t$]\label{thm:multt}
For $f\in\Sw$,
\begin{equation}
\F[tf](\omega)=i\,\frac{\dd}{\dd\omega}\fh(\omega).
\label{eq:multt}
\end{equation}
\end{theorem}

\begin{proof}
By Lemma~\ref{lem:diffint},
$\frac{\dd}{\dd\omega}\fh=-i\int_{\R}tf(t)e^{-i\omega t}\dd t=-i\F[tf](\omega)$. Multiplying by
$i$ gives \eqref{eq:multt}.
\end{proof}

\begin{theorem}[differentiation in $t$]\label{thm:derivt}
For $f\in\Sw$,
\begin{equation}
\F[f'](\omega)=i\omega\,\fh(\omega).
\label{eq:derivt}
\end{equation}
\end{theorem}

\begin{proof}
Integrating by parts,
\begin{equation}
\int_{\R}f'(t)e^{-i\omega t}\dd t
=\big[f(t)e^{-i\omega t}\big]_{-\infty}^{\infty}
+i\omega\int_{\R}f(t)e^{-i\omega t}\dd t .
\end{equation}
The boundary term vanishes because $f$ decays faster than any power and $|e^{-i\omega t}|=1$.
\end{proof}

\begin{remark}[the two rules point in opposite directions]
\eqref{eq:multt} multiplies in time and differentiates in frequency; \eqref{eq:derivt}
differentiates in time and multiplies in frequency. Confusing them is the most common error in this
subject. A quick check that catches the confusion: applying the transform to a Gaussian should not
change the Gaussian's decay rate in a way that contradicts either identity.
\end{remark}

\section{From the constant to every polynomial}

\begin{theorem}\label{thm:tn}
For $n=0,1,2,\dots$,
\begin{equation}
\F[t^{\,n}]=2\pi i^{\,n}\delta^{(n)}(\omega).
\label{eq:tn}
\end{equation}
Consequently, for a polynomial $p(t)=\sum_{n=0}^{N}a_{n}t^{n}$,
\begin{equation}
\F[p](\omega)=2\pi\sum_{n=0}^{N}a_{n}i^{\,n}\delta^{(n)}(\omega).
\label{eq:poly}
\end{equation}
\end{theorem}

\begin{proof}
Induction on $n$. For $n=0$ the statement is $\F[1]=2\pi\delta$, which is \eqref{eq:F1}. Suppose
the formula holds for $n$. Then $t^{n+1}=t\cdot t^{n}$, and by \eqref{eq:multt} applied
(distributionally, see \eqref{eq:multt-dist}) to $t^{n}$,
\begin{equation}
\F[t^{n+1}]=i\frac{\dd}{\dd\omega}\F[t^{n}]
=i\frac{\dd}{\dd\omega}\big(2\pi i^{\,n}\delta^{(n)}\big)
=2\pi i^{\,n+1}\delta^{(n+1)} .
\end{equation}
Linearity gives \eqref{eq:poly}.
\end{proof}

The first few cases, written out, are worth having in view:
\begin{equation}
\F[t]=2\pi i\,\delta',\qquad
\F[t^{2}]=-2\pi\,\delta'',\qquad
\F[t^{3}]=-2\pi i\,\delta''',\qquad
\F[t^{4}]=2\pi\,\delta^{(4)} .
\end{equation}
The alternation of the coefficients is exactly the alternation of $i^{\,n}$.

\begin{remark}[what has not been done]
Equation \eqref{eq:tn} is not obtained by substituting a polynomial into a convergent integral;
\eqref{eq:t-truncated} shows that this would fail. It is obtained by an algebraic rule, whose
legitimacy rests on the definition of $\F$ on distributions (Chapter~\ref{ch:tespered}). Until that
chapter, \eqref{eq:tn} is a computation consistent with all the evidence available, and it is flagged
as such.
\end{remark}

\section{Support, and what lives at a single point}

\begin{definition}[support]\label{def:support}
A distribution $T$ \emph{vanishes on an open set} $U$ if $\ip{T}{\varphi}=0$ for every
$\varphi\in C_{c}^{\infty}$ with $\supp\varphi\subset U$. The \emph{support} $\supp T$ is the
complement of the largest open set on which $T$ vanishes.
\end{definition}

\begin{proposition}\label{prop:suppdelta}
$\supp\delta^{(n)}=\{0\}$ for every $n\ge0$.
\end{proposition}

\begin{proof}
If $\varphi$ vanishes in a neighbourhood of $0$, then $\varphi^{(n)}(0)=0$, and by \eqref{eq:deltan}
$\ip{\delta^{(n)}}{\varphi}=0$. Hence $\delta^{(n)}$ vanishes on $\R\setminus\{0\}$, so its support
is contained in $\{0\}$. It is not empty, since $\ip{\delta^{(n)}}{\varphi}\neq0$ for a suitable
$\varphi$ supported near $0$.
\end{proof}

The converse is more informative than the proposition itself: a distribution supported at a single
point must be a finite combination of $\delta$ and its derivatives. This is the precise content of
the phrase ``a polynomial has spectrum concentrated at zero frequency''.

\begin{theorem}[structure of point-supported distributions]\label{thm:support0}
If $T$ is a distribution with $\supp T\subseteq\{0\}$, then there exist an integer $N\ge0$ and
constants $c_{0},\dots,c_{N}$ such that
\begin{equation}
T=\sum_{n=0}^{N}c_{n}\delta^{(n)} .
\end{equation}
\end{theorem}

\begin{proof}
Let $m$ be the order of $T$ near the origin: by the continuity of $T$ there are $C<\infty$ and
$m\in\mathbb{N}$ with
\begin{equation}
|\ip{T}{\psi}|\le C\max_{0\le j\le m}\|\psi^{(j)}\|_{\infty}
\qquad\text{whenever }\supp\psi\subset[-2,2].
\label{eq:order}
\end{equation}
Choose $N\ge m$ and $\chi\in C_{c}^{\infty}$ with $\chi\equiv1$ on $[-1,1]$. For
$\varphi\in C_{c}^{\infty}$, Taylor's formula with integral remainder gives
\begin{equation}
\varphi(x)=\sum_{n=0}^{N}\frac{\varphi^{(n)}(0)}{n!}x^{n}+x^{N+1}\psi(x),
\qquad
\psi(x)=\int_{0}^{1}\frac{(1-s)^{N}}{N!}\varphi^{(N+1)}(sx)\dd s ,
\end{equation}
and $\psi$ is smooth. Multiply by $\chi$ and split:
\begin{equation}
\varphi=\underbrace{\chi\sum_{n=0}^{N}\frac{\varphi^{(n)}(0)}{n!}x^{n}}_{=:P}
+\underbrace{x^{N+1}\psi\chi}_{=:R}.
\end{equation}
Now $P$ is a polynomial vanishing outside $[-2,2]$ after multiplication by $\chi$, and $R$ vanishes
to order $N+1\ge m+1$ at the origin, so by \eqref{eq:order} $\ip{T}{R}=0$. Therefore
\begin{equation}
\ip{T}{\varphi}=\ip{T}{P}
=\sum_{n=0}^{N}\frac{\varphi^{(n)}(0)}{n!}\ip{T}{\chi x^{n}}
=\sum_{n=0}^{N}\tilde c_{n}\,\varphi^{(n)}(0)
=\sum_{n=0}^{N}\big((-1)^{n}\tilde c_{n}\big)\ip{\delta^{(n)}}{\varphi},
\end{equation}
using \eqref{eq:deltan}. Setting $c_{n}=(-1)^{n}\tilde c_{n}$ finishes the proof.
\end{proof}

\begin{corollary}\label{cor:support-poly}
The spectrum of a polynomial is supported at $\{0\}$, and it is a finite combination of
$\delta^{(n)}$; Theorem~\ref{thm:tn} exhibits the coefficients.
\end{corollary}

\begin{remark}[a check by differential equation, and its limits]
If $p$ has degree $N$ then $p^{(N+1)}=0$, so by \eqref{eq:derivt} in its distributional form,
\begin{equation}
(i\omega)^{N+1}\widehat p=0 ,
\qquad\text{i.e.}\qquad
\omega^{N+1}\widehat p=0 .
\end{equation}
Proposition~\ref{prop:keromega} below then shows that $\widehat p$ is a combination of
$\delta,\delta',\dots,\delta^{(N)}$. This recovers the support statement but \emph{not} the
coefficients: the differential equation is compatible with any choice of the $c_{n}$. The
coefficients come from the recursion \eqref{eq:tn}, and only from it.
\end{remark}

\begin{proposition}\label{prop:keromega}
Let $T$ be a distribution with $\omega T=0$, where $\omega T$ denotes multiplication by the
function $\omega\mapsto\omega$. Then $T=c\,\delta$ for some constant $c$.
\end{proposition}

\begin{proof}
Choose $\chi\in C_{c}^{\infty}$ with $\chi\equiv1$ near $0$. For $\varphi\in C_{c}^{\infty}$ define
\begin{equation}
\psi(\omega)=\frac{\varphi(\omega)-\varphi(0)\chi(\omega)}{\omega},
\end{equation}
which is smooth (the numerator vanishes at $0$ together with its first derivative) and compactly
supported. Hence $\varphi=\varphi(0)\chi+\omega\psi$ and
\begin{equation}
\ip{T}{\varphi}=\varphi(0)\ip{T}{\chi}+\ip{\omega T}{\psi}
=c\,\varphi(0)+0=c\ip{\delta}{\varphi},
\end{equation}
with $c=\ip{T}{\chi}$.
\end{proof}

\begin{remark}
The same argument shows why nothing larger can occur: $\omega\delta'=-\delta\neq0$, so a term
$\delta'$ would destroy the relation $\omega T=0$. Only $\delta$ itself is annihilated by
multiplication by $\omega$.
\end{remark}

\section{Smooth approximations to \texorpdfstring{$\delta'$}{delta prime}}

It is useful to see $\delta'$ as a limit of ordinary functions, because the picture explains both
the sign and the failure of ``area'' as a descriptive tool.

Differentiate the Cauchy kernel $K_{\ep}$ of \eqref{eq:Keps}:
\begin{equation}
K_{\ep}'(x)=\frac{1}{\pi}\cdot\frac{-2\ep x}{(x^{2}+\ep^{2})^{2}}
=-\frac{2\ep x}{\pi\,(x^{2}+\ep^{2})^{2}} .
\label{eq:Kprime}
\end{equation}

\begin{proposition}\label{prop:Kprime}
For every $\ep>0$:
\begin{enumerate}
\item $K_{\ep}'$ vanishes only at $0$, is positive for $x<0$ and negative for $x>0$;
\item its extrema occur at $x=\pm\ep/\sqrt3$, where the value is
$\dfrac{9}{8\sqrt3\,\pi\ep^{2}}$;
\item $\displaystyle\int_{\R}K_{\ep}'(x)\dd x=0$;
\item $\displaystyle\int_{\R}|K_{\ep}'(x)|\dd x=\frac{2}{\pi\ep}\to\infty$;
\item for every bounded continuous $\varphi$,
$\ip{K_{\ep}'}{\varphi}=-\ip{K_{\ep}}{\varphi'}\to-\varphi'(0)=\ip{\delta'}{\varphi}$.
\end{enumerate}
\end{proposition}

\begin{proof}
(i) is clear from \eqref{eq:Kprime}. (ii) Differentiating the logarithmic derivative of
$-x/(x^{2}+\ep^{2})^{2}$ gives the condition $\ep^{2}-3x^{2}=0$; substituting yields
$9/(8\sqrt3\pi\ep^{2})$. (iii) $K_{\ep}'$ is the derivative of a function vanishing at $\pm\infty$,
so its integral is $K_{\ep}(\infty)-K_{\ep}(-\infty)=0$. (iv) By symmetry and the substitution
$u=x/\ep$,
\begin{equation}
\int_{\R}|K_{\ep}'|\dd x
=2\int_{0}^{\infty}\frac{2\ep x}{\pi(x^{2}+\ep^{2})^{2}}\dd x
=\frac{4\ep}{\pi}\cdot\frac{1}{2\ep^{2}}
=\frac{2}{\pi\ep}.
\end{equation}
(v) Integration by parts, justified because $K_{\ep}$ is smooth and decays, gives
$\ip{K_{\ep}'}{\varphi}=-\ip{K_{\ep}}{\varphi'}$, and $\varphi'\in C_{b}$ is bounded and continuous,
so Theorem~\ref{thm:sampling} applies to $\varphi'$.
\end{proof}

Part (iii) is the point. A nonnegative kernel representing $\delta$ is characterized by unit mass
concentrating at the origin; there is no analogous description of $\delta'$, because the total
signed mass is zero identically. What survives is the \emph{shape} of the two lobes, and the shape
encodes exactly one number: the derivative of the test function at the origin.

Part (v) also shows that differentiation is compatible with distributional limits: differentiating a
convergent family produces a convergent family. This will be stated generally in
Proposition~\ref{prop:derivcont}.

\section{Worked examples}

\begin{example}\label{ex:poly1}
Compute $\F[3-2t+t^{2}]$. By linearity and Theorem~\ref{thm:tn},
\begin{equation}
\F[3]-2\F[t]+\F[t^{2}]
=3\cdot2\pi\delta-2\cdot 2\pi i\,\delta'-2\pi\,\delta''
=6\pi\delta-4\pi i\,\delta'-2\pi\,\delta'' .
\end{equation}
Consistency check: $3+t^{2}$ is real and even, and contributes the real even part
$6\pi\delta-2\pi\delta''$; $-2t$ is real and odd and contributes the purely imaginary odd part
$-4\pi i\delta'$. This is the symmetry statement of Proposition~\ref{prop:symmetries}, in
distributional form.
\end{example}

\begin{example}\label{ex:delta-prime-test}
Let $\varphi(x)=e^{-x^{2}}+xe^{-x^{2}}$. Then
\begin{equation}
\varphi(0)=1,\qquad
\varphi'(x)=e^{-x^{2}}\big(1-2x-2x^{2}\big),\qquad
\varphi'(0)=1,\qquad
\varphi''(0)=-2 .
\end{equation}
Hence
\begin{equation}
\ip{\delta}{\varphi}=1,\qquad
\ip{\delta'}{\varphi}=-1,\qquad
\ip{\delta''}{\varphi}=-2 .
\end{equation}
\end{example}

\begin{example}\label{ex:ladder}
The identities
\begin{equation}
x\,\delta=0,\qquad
x\,\delta'=-\delta,\qquad
x\,\delta^{(n)}=-n\,\delta^{(n-1)}
\label{eq:xdeltan}
\end{equation}
are proved in Chapter~\ref{ch:dist}. Granting them for a moment, they explain the structure found in
Proposition~\ref{prop:keromega}: multiplication by $\omega$ lowers the order of $\delta^{(n)}$ by
one, so only $\delta^{(0)}$ survives annihilation.
\end{example}

\section{Summary of the chapter}

\begin{enumerate}
\item Polynomials cannot be substituted into \eqref{eq:ftpair}; the truncation grows like $T^{n}$.
\item Differentiation of distributions is defined by transferring the derivative to the test
function, at the price of $(-1)^{n}$ (Definition~\ref{def:distderiv}).
\item Multiplication by $t$ in time corresponds to differentiation in frequency
(Theorem~\ref{thm:multt}); differentiation in time corresponds to multiplication by $i\omega$
(Theorem~\ref{thm:derivt}).
\item The recursion gives $\F[t^{n}]=2\pi i^{n}\delta^{(n)}$
(Theorem~\ref{thm:tn}), and linearity extends it to all polynomials.
\item A distribution supported at the origin is a finite combination of $\delta^{(n)}$
(Theorem~\ref{thm:support0}); this is what ``support at zero frequency'' means.
\item Differentiating the approximating kernel yields $\delta'$ as a distributional limit with two
lobes, zero total mass, and divergent total variation
(Proposition~\ref{prop:Kprime}).
\end{enumerate}
